\scanpage{200}
\mspage{9.}
au préfaisceau $S''\longrightarrow\mathrm{Aut}\,x_{S''}$ est représentable par un \deleted{objet groupe de la catégorie des} rev. étale de $S'$. On a une ``représentation'' de cette gerbe dans le préchamp $\mathrm{Fl}_S$. Á ces \unsure{belles|telles} données, on associe une donnée de ramification sur $S$ ainsi

Les objets de $\mathscr R_S$ sont les hom. cartésiens \uline{sur $\mathrm{Fl}_S$} $\mathscr G\longrightarrow\mathrm{Rev}\,\mathrm{Fl}_S$ qui relèvent $\mathscr G\unsure{-|\to}\mathrm{Fl}_S$, et itou quand $r$ est remplacé par $r'$.
\diagram{\begin{tikzcd}[ampersand replacement=\&]
{}\&{}\&\mathrm{Rev}\,\mathrm{Fl}_S\arrow[d]\\
\mathscr G\arrow[rr]\arrow[dr,dash]\&{}\&\mathrm{Fl}_S\arrow[dl]\\
{}\&Et(S)\&{}
\end{tikzcd}}
Le foncteur $r_S:\mathscr R_S\longrightarrow (Sch)_{/S}$ est le foncteur $\varinjlim$. On vérifie que tt marche.

\uline{Ex 1} Donnons \deleted{\unsure{la}} famille $\underline a=(a_i)_{i\in\unsure{I|J}}$ de sections de $\mathcal O_S$, et famille $\underline n=(\underline n_i)_{i\in I}$ d'entiers $\geq1$, \inserted{premiers aux car. résiduelles} \deleted{\unsure{les points}} et tous égaux à 1 sauf un nb fini. \deleted{\unsure{tout sauf}} Soit $C_{S'}=\deleted{\unsure{F\!l}}$ groupoïde défini par le groupe $\mu_{\underline n}(S')=\prod_{i\in I}\mu_{n_i}(S)$, \deleted{Supposons}\inserted{(unique objet $e_{S'}$)} et $s(e_{S'})=Z_{\underline n}^{\underline a}$, avec l'opération qu'on devine de ce groupe. Les axiomes sont vérifiés. [En fait, cela correspond au cas d'une \uline{gerbe triviale} définie par $Z_{\underline n}^{\underline a}$ sur lequel $\mu_{\underline n}^{\underline a}$ opère ].

\uline{Ex 1 bis} \inserted{Donnée \deleted{de} de ram. Kummer définie par $\underline a,\underline n$.} Comme avant, mais on passe à la limite sur $\underline n$ variable, $\underline a$ fixé ; donnée de ramification kummérienne définie par $\underline a$.

\uline{Ex} 2 Donnons famille $(D_i)_{i\in I}$ de diviseurs, et famille $(n_i)_{i\in I}$ d'entiers satisfaisant mêmes conditions que dans ex. \uline{1}. On prend
% archive sheet number: 199
